\section{\minerva Experiment}
\label{sec:expt}
The \minerva experiment combines a fine-grained tracking detector with the high-intensity 
NuMI beam line~\cite{Anderson:1998zza} and the MINOS near detector~\cite{Michael:2008bc}.  
The neutrino beam is created by directing 120 GeV protons onto a graphite target, producing charged 
particles (mostly pions and kaons) which are focused into a 
beam by two magnetic horns.  Downstream of the horns, most of the pions and kaons decay within 
the 675 m helium-filled decay pipe to produce neutrinos.  Approximately 97\% of the muon neutrinos that enter \minerva 
are produced by pion decay, with the remainder produced from kaon decay.

The \minerva detector consists of a central tracking volume 
preceded by nuclear targets, which are not used in this analysis, and surrounded by electromagnetic 
and hadronic calorimeters.
In the tracking volume, triangular polystyrene scintillator strips with a 1.7~cm strip-to-strip pitch are arranged into planes
arrayed perpendicularly to the horizontal axis, which is inclined by 3.4$^\circ$
relative to the beam direction.  
Three plane orientations, at $0^\circ$ and $\pm 60^\circ$ relative to the vertical axis, enable 
unambiguous 3-dimensional reconstruction of the neutrino
interaction point and charged particle tracks.  Each scintillator strip contains a wavelength-shifting fiber 
that is read out by a multi-anode photomultiplier tube.  The 3.0 ns timing resolution of the readout electronics is 
adequate for separating multiple interactions within a single beam spill.   
The MINOS near detector, located 2~m downstream of the \minerva detector, 
is used to reconstruct muon momentum and charge.  More information on the design, calibration, and performance of 
the \minerva detector, including the elemental composition of the tracking volume, is provided in Ref.~\cite{minerva_nim}.

The data for this measurement were collected 
between March 2010 and April 2012 and correspond to an integrated $3.04\times 10^{20}$ protons on 
target (POT).  For these data the horn current was configured to produce 
a muon neutrino beam, and the MINOS detector's magnet polarity was set to focus negative muons.   

\section{Experiment Simulation}
\label{sec:expt_sim}
The neutrino beam is simulated by a GEANT4-based model~\cite{Agostinelli2003250,1610988} that is tuned to agree with 
hadron production measurements on carbon~\cite{Alt:2006fr, Lebedev:2007zz} by the procedure described in Ref.~\cite{MINERVAnubarprl}.  Uncertainty on
the neutrino flux is determined by the precision in these measurements, uncertainties in the beam line
focusing system and alignment~\cite{Pavlovic:2008zz}, and comparisons between different hadron production models in regions not
covered by the hadron production data referenced above.  The integrated neutrino flux over the range 
$1.5\le E_\nu \le 10$~GeV is estimated to be $\unit[2.77 \times 10^{-8}]{cm^{-2}/POT}$.  Table~\ref{tab:nu_flux} lists the flux as a function of energy.

\begingroup
\squeezetable
\begin{table*}[t]
\centering
\begin{tabular}{l|ccccccccccccc}
\hline \hline
$E_\nu$ (GeV) & 
$1.5 - 2$ &
$2 - 2.5$ &
$2.5 - 3$ &
$3 - 3.5$ &
$3.5 - 4$ &
$4 - 4.5$ &
$4.5 - 5$ &
$5 - 5.5$ \\
Flux ($\nu_{\mu}$/cm$^2$/POT ($\times 10^{-8}$)) &
$0.291	$ &
$0.387	$&
$0.476	$&
$0.502   $&
$0.402	$&
$0.242	$&
$0.131	$&
$0.077	$ \\
\hline
$E_\nu$ (GeV) & 
$5.5 - 6$ &
$6 - 6.5$ &
$6.5 - 7$ &
$7 - 7.5$ &
$7.5 - 8$ &
$8 - 8.5$ &
$8.5 - 9$ &
$9 - 9.5$ &
$9.5 - 10$ \\
Flux ($\nu_{\mu}$/cm$^2$/POT ($\times 10^{-8}$)) &
$0.053   $&
$0.041	$&
$0.035	$&
$0.030	$&
$0.026	$&
$0.023	$&
$0.021	$&
$0.019	$&
$0.017	$ \\ \hline \hline
\end{tabular}
\caption{The $\nu_\mu$ flux per POT for the data included in this analysis.}
\label{tab:nu_flux}
\end{table*}
\endgroup

The \minerva detector response is also simulated by a GEANT4-based model.  The muon energy loss scale of the detector
is known to within 2\% by requiring agreement between data and simulation of both the photon statistics and the reconstructed
energy deposited by momentum-analyzed throughgoing muons. Calorimetric corrections used to reconstruct the
energy of hadronic showers are determined from the simulation by the procedure described in Ref.~\cite{minerva_nim}.  The 
uncertainties on the hadron interaction models in GEANT4 are determined to be $\sim$10\% by external 
data~\cite{Ashery:1981tq,Allardyce:1973ce,Saunders:1996ic,Lee:2002eq}.  The tracking efficiency and 
energy response of single hadrons, as well as the scintillation Birks constant, are determined from measurements made 
with a scaled-down replica of the MINERvA detector 
in a low energy hadron test beam~\cite{mnv_testbeam}. The response of the MINOS near detector to muons is determined by a tuned
GEANT-based
simulation~\cite{Michael:2008bc}.

Neutrino interactions are simulated using the GENIE 2.6.2 neutrino event generator~\cite{Andreopoulos201087}.  Non-coherent interactions 
are treated as neutrino-nucleon scattering within a relativistic Fermi gas.  The nucleon momentum distribution is 
modified with a high-energy tail to account for nucleon-nucleon interactions, but interactions with correlated 
nucleon pairs are not included in the simulation.  Pauli-blocking is applied to quasielastic and elastic scattering, but not to 
resonance baryon production.  The structure functions in the deep inelastic scattering (DIS) model are modified to reproduce 
the shadowing, anti-shadowing, and EMC effects observed in charged-lepton nuclear scattering data.

Almost all pion-production events observed in \minerva are due to baryon resonance production, nonresonant pion production (including DIS), 
and coherent pion production.  For baryon resonance production at $W<1.7$~GeV, the formalism of 
Rein-Sehgal~\cite{Rein:1980wg} is used with modern resonance properties~\cite{pdg} and an axial mass $M_{A}=1.12\pm0.22$~GeV.  
However, GENIE differs from Rein-Sehgal in a couple ways. Resonance interference and lepton mass terms in the 
cross section calculation are not included. 
Most signficantly, the angular spectrum of the $\Delta$ decay is nominally isotropic in GENIE; this analysis instead reweights GENIE such 
that the $\Delta$ decay angular anisotropy is half that predicted by 
Rein-Sehgal.  Excursions from isotropic to the full Rein-Sehgal anisotropic prediction are included as a systematic uncertainty.
Nonresonant pion production is simulated using the Bodek-Yang 
model~\cite {Bodek:2004pc} and is constrained below $W=1.7$~GeV by 
neutrino-deuterium bubble chamber data. Coherent pion production is described according to the 
model of Rein-Sehgal modified with lepton mass terms~\cite{Rein:2007}.  Uncertainties on the components of the neutrino-interaction model are 
provided by GENIE. 

Pion and nucleon FSI processes are modeled in GENIE using an effective
intranuclear cascade model~\cite{Dytman:2011zz}, called the ``hA'' model, that simulates the full
cascade as a single interaction and tunes the overall interaction rate
to hadron-nucleus total reaction cross section data.  For 
light nuclei such as carbon, a single interaction happens for a large
fraction of the events.  The final state particle multiplicity and kinematic 
distributions are also tuned to data.  
This model has good agreement with a wide range 
of data from hadron-nucleus scattering experiments for many 
targets. 
Uncertainties in the FSI model are evaluated by varying its parameters
within measured uncertainties~\cite{Lee:2002eq,Ashery:1981tq}.


 




